% =====================================================================
\chapter{Ángulos, razones trigonométricas y circunferencia unidad}
\label{tema:1}
% =====================================================================

Este tema fija la medida de ángulos en radianes, define las razones trigonométricas primero en el triángulo rectángulo y después, para cualquier número real, mediante la circunferencia unidad, y deduce de las simetrías de la circunferencia las propiedades básicas de signo, paridad y periodicidad. Se sigue el capítulo~7 de~\cite{osat} (§7.1 a §7.4) y el doble enfoque, del triángulo y de la circunferencia unidad, de Stewart, Redlin y Watson~\cite{stewart}.

\section{Ángulos y radianes}

Un ángulo en posición estándar tiene el vértice en el origen y el lado inicial sobre el semieje positivo de abscisas; es positivo si se recorre en sentido antihorario y negativo en sentido horario. Dos ángulos son \emph{coterminales} si tienen el mismo lado final.

\begin{definicion}[Radián]\label{def:radian}
Sea una circunferencia de radio $r>0$ y un ángulo central que abarca un arco de longitud $s$. La medida del ángulo en \emph{radianes} es $\theta=s/r$. Una vuelta completa mide $2\pi$ radianes, de modo que
\[
180^{\circ}=\pi\ \text{rad},\qquad 1^{\circ}=\frac{\pi}{180}\ \text{rad}\approx 0{,}017\,45\ \text{rad},\qquad 1\ \text{rad}=\frac{180^{\circ}}{\pi}\approx 57{,}30^{\circ}.
\]
\end{definicion}

La definición no depende del radio elegido: al multiplicar $r$ por $\lambda$, la semejanza multiplica también $s$ por $\lambda$, y el cociente no cambia. El radián es, por tanto, una razón entre dos longitudes y carece de dimensiones.

\begin{proposicion}[Arco y sector]\label{prop:arco-sector}
En una circunferencia de radio $r$, un ángulo central de $\theta$ radianes ($0\le\theta\le 2\pi$) abarca un arco de longitud $s=r\theta$ y un sector circular de área $A=\tfrac12 r^{2}\theta$.
\end{proposicion}

\begin{proof}
La primera igualdad es la \cref{def:radian} despejada. Para la segunda, el área del sector es proporcional al ángulo central; el círculo completo ($\theta=2\pi$) tiene área $\pi r^{2}$, luego $A=\pi r^{2}\cdot\theta/(2\pi)=\tfrac12 r^{2}\theta$.
\end{proof}

\begin{ejemplo}\label{ej:arco}
En una circunferencia de radio $r=4$, el ángulo de $45^{\circ}=\pi/4$ rad abarca un arco $s=4\cdot\pi/4=\pi\approx 3{,}142$ y un sector de área $\tfrac12\cdot 16\cdot\pi/4=2\pi\approx 6{,}283$.
\end{ejemplo}

\section{Razones trigonométricas en el triángulo rectángulo}

\begin{definicion}[Razones de un ángulo agudo]\label{def:razones-triangulo}
Sea un triángulo rectángulo con un ángulo agudo $\theta$, cateto opuesto $a$, cateto contiguo $b$ e hipotenusa $c$. Se definen
\[
\sin\theta=\frac{a}{c},\qquad \cos\theta=\frac{b}{c},\qquad \tan\theta=\frac{a}{b},
\]
y sus inversas aritméticas $\csc\theta=1/\sin\theta$, $\sec\theta=1/\cos\theta$, $\cot\theta=1/\tan\theta$.
\end{definicion}

\begin{proposicion}\label{prop:razones-bien-definidas}
Las razones de la \cref{def:razones-triangulo} solo dependen de $\theta$, no del triángulo rectángulo elegido.
\end{proposicion}

\begin{proof}
Dos triángulos rectángulos con un ángulo agudo $\theta$ común tienen sus tres ángulos iguales ($\pi/2$, $\theta$ y $\pi/2-\theta$), luego son semejantes. En triángulos semejantes los lados correspondientes son proporcionales, $a'=\lambda a$, $b'=\lambda b$, $c'=\lambda c$, y los cocientes no cambian: $a'/c'=a/c$, etc.
\end{proof}

\begin{proposicion}[Valores notables]\label{prop:notables}
Se cumplen $\sin\frac{\pi}{6}=\cos\frac{\pi}{3}=\frac12$, $\sin\frac{\pi}{3}=\cos\frac{\pi}{6}=\frac{\sqrt3}{2}$ y $\sin\frac{\pi}{4}=\cos\frac{\pi}{4}=\frac{\sqrt2}{2}$.
\end{proposicion}

\begin{proof}
Para $\pi/4$: un triángulo rectángulo isósceles de catetos~1 tiene, por Pitágoras, hipotenusa $\sqrt2$, luego $\sin\frac{\pi}{4}=\cos\frac{\pi}{4}=1/\sqrt2=\sqrt2/2$. Para $\pi/6$ y $\pi/3$: la altura de un triángulo equilátero de lado~2 lo divide en dos triángulos rectángulos de hipotenusa~2, cateto~1 (opuesto al ángulo de $\pi/6$) y otro cateto $\sqrt{4-1}=\sqrt3$ (opuesto al ángulo de $\pi/3$). Las razones se leen directamente.
\end{proof}

\section{La circunferencia unidad}

La \cref{def:razones-triangulo} solo tiene sentido para $0<\theta<\pi/2$. Para extenderla a cualquier número real se usa la circunferencia de radio~1 centrada en el origen, $x^{2}+y^{2}=1$ (\cref{fig:circunferencia}).

\begin{definicion}[Seno y coseno de un número real]\label{def:circunferencia}
Para $\theta\in\R$, sea $P(\theta)$ el punto de la circunferencia unidad al que se llega recorriendo sobre ella, desde $(1,0)$, un arco de longitud $\abs{\theta}$, en sentido antihorario si $\theta\ge 0$ y horario si $\theta<0$. Se definen $\cos\theta$ y $\sin\theta$ como las coordenadas de ese punto:
\[
P(\theta)=(\cos\theta,\ \sin\theta).
\]
Cuando $\cos\theta\neq 0$ se define $\tan\theta=\sin\theta/\cos\theta$, y análogamente $\sec$, $\csc$ y $\cot$.
\end{definicion}

\begin{figure}[ht]
\centering
\begin{tikzpicture}[scale=2.2,font=\small]
  \draw[->] (-1.3,0) -- (1.35,0) node[right] {$x$};
  \draw[->] (0,-1.2) -- (0,1.3) node[above] {$y$};
  \draw[thick] (0,0) circle (1);
  \coordinate (P) at (40:1);
  \draw[thick,->] (0,0) -- (P);
  \draw[dashed] (P) -- (P|-0,0);
  \draw[dashed] (P) -- (P-|0,0);
  \fill (P) circle (0.025) node[above right] {$P(\theta)=(\cos\theta,\sin\theta)$};
  \draw[very thick,->] (1,0) arc (0:40:1);
  \node at (0.95,0.33) [right] {$\theta$};
  \draw (0.25,0) arc (0:40:0.25);
  \node at (0.38,0.14) {$\theta$};
  \node[below] at ($(P|-0,0)$) {$\cos\theta$};
  \node[left] at ($(P-|0,0)$) {$\sin\theta$};
  \node[below left] at (1,0) {$1$};
  \node at (0.7,1.05) {I};
  \node at (-0.7,1.05) {II};
  \node at (-0.7,-1.05) {III};
  \node at (0.7,-1.05) {IV};
\end{tikzpicture}
\caption{La circunferencia unidad. El arco de longitud $\theta$ desde $(1,0)$ llega al punto $P(\theta)$; sus coordenadas son $\cos\theta$ y $\sin\theta$. Los números romanos indican los cuadrantes.}
\label{fig:circunferencia}
\end{figure}

\begin{proposicion}\label{prop:coherencia}
Para $0<\theta<\pi/2$, la \cref{def:circunferencia} coincide con la \cref{def:razones-triangulo}.
\end{proposicion}

\begin{proof}
Si $0<\theta<\pi/2$, el punto $P(\theta)=(x,y)$ está en el primer cuadrante, y el triángulo de vértices $O$, $(x,0)$ y $P(\theta)$ es rectángulo en $(x,0)$, con ángulo $\theta$ en $O$ (por la \cref{def:radian} con $r=1$), hipotenusa $1$, cateto contiguo $x$ y cateto opuesto $y$. La \cref{def:razones-triangulo} da $\cos\theta=x/1$ y $\sin\theta=y/1$.
\end{proof}

\begin{teorema}[Propiedades básicas]\label{teo:simetrias}
Para todo $\theta\in\R$ y todo $k\in\Z$:
\begin{enumerate}
  \item (identidad pitagórica) $\sin^{2}\theta+\cos^{2}\theta=1$; en particular, $\abs{\sin\theta}\le 1$ y $\abs{\cos\theta}\le 1$;
  \item (periodicidad) $\sin(\theta+2k\pi)=\sin\theta$ y $\cos(\theta+2k\pi)=\cos\theta$; más aún, $P(\theta)=P(\varphi)$ si y solo si $\theta-\varphi\in2\pi\Z$;
  \item (paridad) $\cos(-\theta)=\cos\theta$ y $\sin(-\theta)=-\sin\theta$;
  \item (suplementario) $\sin(\pi-\theta)=\sin\theta$ y $\cos(\pi-\theta)=-\cos\theta$;
  \item (media vuelta) $\sin(\theta+\pi)=-\sin\theta$ y $\cos(\theta+\pi)=-\cos\theta$;
  \item (complementario) $\sin(\tfrac{\pi}{2}-\theta)=\cos\theta$ y $\cos(\tfrac{\pi}{2}-\theta)=\sin\theta$.
\end{enumerate}
\end{teorema}

\begin{proof}
(i) $P(\theta)$ está en la circunferencia $x^{2}+y^{2}=1$. (ii) La circunferencia mide $2\pi$: recorrer $2k\pi$ más da $\abs{k}$ vueltas completas y devuelve al mismo punto. Recíprocamente, todo $\theta$ se escribe como $\theta=\theta_0+2k\pi$ con $\theta_0\in[0,2\pi)$ y $P(\theta)=P(\theta_0)$; como la circunferencia mide exactamente $2\pi$, arcos antihorarios distintos de longitud menor que $2\pi$ desde $(1,0)$ terminan en puntos distintos, luego $P(\theta)=P(\varphi)$ obliga a $\theta_0=\varphi_0$, es decir, $\theta-\varphi\in2\pi\Z$. Las restantes se deducen de que las simetrías del plano que fijan la circunferencia y el punto de partida $(1,0)$, o lo llevan a otro punto conocido, conservan las longitudes de arco. (iii) La reflexión en el eje $x$ fija $(1,0)$ e invierte el sentido de recorrido, luego lleva $P(\theta)$ a $P(-\theta)$: $(x,y)\mapsto(x,-y)$. (iv) La reflexión en el eje $y$, $(x,y)\mapsto(-x,y)$, lleva $(1,0)$ a $(-1,0)=P(\pi)$ e invierte el sentido; un arco de longitud $\theta$ desde $(1,0)$ se transforma en uno de longitud $\theta$ desde $P(\pi)$ en sentido horario, que termina en $P(\pi-\theta)$. (v) El giro de $\pi$ alrededor del origen, $(x,y)\mapsto(-x,-y)$, lleva $(1,0)$ a $P(\pi)$ y conserva el sentido; por tanto lleva $P(\theta)$ a $P(\theta+\pi)$. (vi) La reflexión en la recta $y=x$, $(x,y)\mapsto(y,x)$, lleva $(1,0)$ a $(0,1)=P(\pi/2)$ e invierte el sentido; lleva $P(\theta)$ a $P(\pi/2-\theta)$, de modo que $(\cos(\frac{\pi}{2}-\theta),\sin(\frac{\pi}{2}-\theta))=(\sin\theta,\cos\theta)$.
\end{proof}

\begin{corolario}[Signos por cuadrantes]\label{cor:signos}
El seno es positivo en los cuadrantes I y II; el coseno, en I y IV; la tangente, en I y III.
\end{corolario}

\begin{proof}
El seno es la ordenada de $P(\theta)$ y el coseno la abscisa; la tangente es su cociente, positiva cuando ambos tienen el mismo signo.
\end{proof}

\begin{definicion}[Ángulo de referencia]
El \emph{ángulo de referencia} de $\theta$ es el ángulo agudo $\theta'\in[0,\pi/2]$ que forma el lado final de $\theta$ con el eje $x$. Por el \cref{teo:simetrias}, $\abs{\sin\theta}=\sin\theta'$ y $\abs{\cos\theta}=\cos\theta'$; el signo lo da el \cref{cor:signos}.
\end{definicion}

\begin{ejemplo}[Valores por reducción]\label{ej:reduccion}
Para $\theta=5\pi/4$: es coterminal consigo mismo, está en el cuadrante III y su ángulo de referencia es $\pi/4$. Por tanto $\cos\frac{5\pi}{4}=\sin\frac{5\pi}{4}=-\frac{\sqrt2}{2}$ y $\tan\frac{5\pi}{4}=1$. Para $\theta=-\pi/6$, por la paridad, $\sin(-\frac{\pi}{6})=-\frac12$ y $\cos(-\frac{\pi}{6})=\frac{\sqrt3}{2}$.
\end{ejemplo}

\begin{table}[ht]
\centering\small
\caption{Valores notables. Los de los demás cuadrantes se obtienen con el \cref{teo:simetrias}.}\label{tab:notables}
\begin{tabular}{@{}lccccc@{}}
\toprule
$\theta$ & $0$ & $\pi/6$ & $\pi/4$ & $\pi/3$ & $\pi/2$ \\
\midrule
$\sin\theta$ & $0$ & $1/2$ & $\sqrt2/2$ & $\sqrt3/2$ & $1$ \\
$\cos\theta$ & $1$ & $\sqrt3/2$ & $\sqrt2/2$ & $1/2$ & $0$ \\
$\tan\theta$ & $0$ & $\sqrt3/3$ & $1$ & $\sqrt3$ & no definida \\
\bottomrule
\end{tabular}
\end{table}

\begin{observacion}[Errores frecuentes]
(a) Usar grados donde se esperan radianes: en todo el curso, y en cualquier lenguaje de programación, $\sin(30)$ es el seno de 30 radianes, que vale $-0{,}988$, no $1/2$. (b) Confundir $\sin^{2}x=(\sin x)^{2}$ con $\sin(x^{2})$. (c) Leer $\sin^{-1}x$ como $1/\sin x$: esa notación designa la función inversa $\arcsin x$ (\cref{tema:4}); el inverso aritmético es $\csc x$. (d) Olvidar el signo del cuadrante al usar el ángulo de referencia.
\end{observacion}

\begin{observacion}[Conexión con el curso]
Los ángulos que aparecen en el curso se miden siempre en radianes. En la codificación rotatoria RoPE (módulo~14, tema~4), el par $i$ de coordenadas de la consulta en la posición $m$ se gira un ángulo $m\theta_i$ con $\theta_i=10000^{-2i/d}$ radianes; para $i=0$ el ángulo crece un radián por posición y da una vuelta completa cada $2\pi\approx 6{,}28$ posiciones. El ángulo entre dos embeddings (módulo~2, tema~2) se obtiene también en radianes, en el intervalo $[0,\pi]$.
\end{observacion}

\begin{lecturas}
\item OpenStax, \emph{Algebra and Trigonometry}~\cite{osat}, §7.1 «Angles», §7.2 «Right Triangle Trigonometry», §7.3 «Unit Circle» y §7.4 «The Other Trigonometric Functions».
\item Stewart, Redlin y Watson~\cite{stewart}, capítulos dedicados a las funciones trigonométricas desde la circunferencia unidad y desde el triángulo rectángulo.
\end{lecturas}

\section*{Ejercicios del tema 1}
\addcontentsline{toc}{section}{Ejercicios del tema 1}

\begin{ejercicio}{1}{eje:1-conversion}
Exprese en radianes los ángulos $135^{\circ}$, $210^{\circ}$, $-36^{\circ}$ y $1^{\circ}$, y en grados los ángulos $5\pi/6$, $7\pi/4$ y $2$ rad.
\end{ejercicio}

\begin{ejercicio}{1}{eje:1-cuadrantes}
Sin calculadora, determine el cuadrante y el signo de $\sin\theta$, $\cos\theta$ y $\tan\theta$ para $\theta=2$, $\theta=4$ y $\theta=5$ (en radianes). Use $\pi\approx 3{,}14$.
\end{ejercicio}

\begin{ejercicio}{1}{eje:1-vf}
Justifique si son verdaderas o falsas: (a) $\sin^{2}x=\sin(x^{2})$ para todo $x$; (b) $\sin^{-1}x=1/\sin x$; (c) existe $\theta$ con $\cos\theta=1{,}2$; (d) si $\theta$ y $\varphi$ son coterminales, $\tan\theta=\tan\varphi$ siempre que estén definidas.
\end{ejercicio}

\begin{ejercicio}{2}{eje:1-arco}
En una circunferencia de radio $6$ cm, calcule la longitud del arco y el área del sector correspondientes a un ángulo central de $150^{\circ}$. ¿Qué ángulo central, en radianes y en grados, abarca un arco de 9 cm?
\fuente{adaptado de OpenStax, \emph{Algebra and Trigonometry}~\cite{osat}, §7.1, ejercicios de la sección sobre longitud de arco y área del sector.}
\end{ejercicio}

\begin{ejercicio}{2}{eje:1-exactos}
Calcule el valor exacto de $\sin\frac{5\pi}{6}$, $\cos\frac{7\pi}{4}$, $\tan\frac{4\pi}{3}$, $\sin(-\frac{\pi}{3})$, $\cos\frac{11\pi}{6}$ y $\sin\frac{3\pi}{2}$, indicando en cada caso el cuadrante y el ángulo de referencia.
\fuente{adaptado de OpenStax, \emph{Algebra and Trigonometry}~\cite{osat}, §7.3, ejercicios de la sección sobre valores en la circunferencia unidad.}
\end{ejercicio}

\begin{ejercicio}{2}{eje:1-seis-razones}
Se sabe que $\sin\theta=3/5$ y que $\theta$ está en el segundo cuadrante. Calcule $\cos\theta$, $\tan\theta$, $\sec\theta$, $\csc\theta$ y $\cot\theta$.
\fuente{adaptado de OpenStax, \emph{Algebra and Trigonometry}~\cite{osat}, §7.4, ejercicios de la sección sobre las demás funciones trigonométricas.}
\end{ejercicio}

\begin{ejercicio}{2}{eje:1-elevacion}
Desde un punto del suelo situado a 30 m de la base de un edificio, la parte alta del edificio se ve con un ángulo de elevación de $60^{\circ}$. Calcule la altura del edificio y la distancia desde el punto de observación hasta la parte alta.
\fuente{adaptado de OpenStax, \emph{Algebra and Trigonometry}~\cite{osat}, §7.2, problemas de aplicación con ángulos de elevación.}
\end{ejercicio}

\begin{ejercicio}{2}{eje:1-coterminales}
Para cada uno de los ángulos $1020^{\circ}$, $-11\pi/4$ y $17\pi/6$, encuentre el ángulo coterminal en $[0^{\circ},360^{\circ})$ o en $[0,2\pi)$, su ángulo de referencia y el valor exacto de su seno y su coseno.
\end{ejercicio}

\begin{ejercicio}{2}{eje:1-rope}
En RoPE, el par de coordenadas más rápido se gira un ángulo de $m$ radianes en la posición $m$. Para $m=100$: (a) calcule cuántas vueltas completas se han dado y el ángulo coterminal en $[0,2\pi)$; (b) indique el cuadrante y el signo del seno y del coseno; (c) ¿cada cuántas posiciones se repite exactamente el ángulo? Justifique la respuesta.
\end{ejercicio}

\begin{ejercicio}{2}{eje:1-simplificar}
Simplifique, para los valores de $x$ en que estén definidas, las expresiones (a) $\sin(\pi-x)+\cos(\tfrac{\pi}{2}+x)$; (b) $\cos(x+\pi)\cos(-x)-\sin(x+\pi)\sin(\pi-x)$; (c) $\tan(\pi+x)\cot(\pi-x)$.
\end{ejercicio}

\begin{ejercicio}{3}{eje:1-sinx-x-tanx}
Demuestre que, para $0<x<\pi/2$, se cumple $\sin x<x<\tan x$. Indicación: compare las áreas de dos triángulos y de un sector con vértice en el origen.
\end{ejercicio}

\begin{ejercicio}{3}{eje:1-abs-sin}
Usando el \cref{eje:1-sinx-x-tanx}, demuestre que $\abs{\sin x}\le\abs{x}$ para todo $x\in\R$, con igualdad solo si $x=0$.
\end{ejercicio}

\begin{ejercicio}{3}{eje:1-sin-mas-cos}
Demuestre que $\abs{\sin\theta+\cos\theta}\le\sqrt2$ para todo $\theta\in\R$, y determine los valores de $\theta$ para los que $\sin\theta+\cos\theta=\sqrt2$.
\end{ejercicio}

\begin{ejercicio}{3}{eje:1-periodo-tan}
(a) Demuestre que $\sin\theta=0$ si y solo si $\theta=k\pi$ con $k\in\Z$. (b) Demuestre que $\tan(\theta+\pi)=\tan\theta$ en todo su dominio y que ningún número $0<T<\pi$ cumple $\tan(\theta+T)=\tan\theta$ para todo $\theta$ del dominio.
\end{ejercicio}

\begin{ejercicio}{3}{eje:1-poligono}
Sea $n\ge 3$. Demuestre que el perímetro del polígono regular de $n$ lados inscrito en la circunferencia unidad es $p_n=2n\sin(\pi/n)$, y que $p_n<2\pi$. Calcule $p_6$ y compruebe el resultado.
\end{ejercicio}
